% Flächen – bank/flaechen/ (zusammenbau v0.4, Heft)
\einheitenkopf[t1]{Flächen}
\setcounter{aufgabe}{0}

% Anlauf (ohne Prüfkennung)
\begin{aufgabe}{Rechteck – Anlauf}
\begin{teile}
\teil Ein Garten bekommt einen Zaun. Fläche oder Umfang? \\ \kreuz{Fläche} \\ \kreuz{Umfang}
\teil Rechteck, $a = 8$ cm, $b = 3$ cm – Fläche und Umfang? A = \leerfeld[cm²] , u = \leerfeld[cm]
\teil Rechteck, $a = 4{,}5$ cm, $b = 2$ cm – Fläche und Umfang? A = \leerfeld[cm²] , u = \leerfeld[cm]
\end{teile}
\end{aufgabe}

\begin{aufgabe}{Rechteck – Prüfungsaufgaben}
\begin{teile}
\teil Ein rechteckiger Bilderrahmen hat einen Umfang von $34$ cm, eine Seite ist $11$ cm lang. Wie lang ist die andere Seite? \hfill (P10 2018 OS) \\ \leerfeld[cm]
\teil Ein rechteckiges Beet ist mit $38$ m Kantenstein eingefasst und $14$ m lang. Wie breit ist das Beet? \hfill (P10 2018 OS) \\ \leerfeld[m]
\end{teile}
\end{aufgabe}

% Anlauf (ohne Prüfkennung)
\begin{aufgabe}{Dreieck – Anlauf}
\begin{teile}
\teil Zeichne die Höhe zur Grundseite $g$ ein und markiere den rechten Winkel.

\dreieck{(0,0)}{(5,0)}{(2,3)}{}{}{$g$}{}{}{}
\teil Dreieck, $g = 8$ cm, $h = 7$ cm – Fläche? A = \leerfeld[cm²]
\teil Rechtwinkliges Dreieck – Fläche?

\dreieckrw{5}{2.67}{$15$ cm}{$8$ cm}{$17$ cm}
A = \leerfeld[cm²]
\end{teile}
\end{aufgabe}

\begin{aufgabe}{Dreieck – Prüfungsaufgaben}
\begin{teile}
\teil Ein dreieckiges Grundstück ABC hat bei C einen rechten Winkel; $BC = 16$ m, $AC = 45$ m. Ein Teilstück QBC mit Q auf AC soll $120$ m² groß sein. Wie lang ist QC? \hfill (P10 2020 OS) \\ QC = \leerfeld[m]
\teil Ein Garten ABC hat bei C einen rechten Winkel; $BC = 9$ m, $AC = 25$ m. Das Dreieck PBC mit P auf AC wird Rasen und soll $63$ m² groß sein. Wie lang ist PC? \hfill (P10 2020 OS) \\ PC = \leerfeld[m]
\end{teile}
\end{aufgabe}

% Anlauf (ohne Prüfkennung)
\begin{aufgabe}{Trapez – Anlauf}
\begin{teile}
\teil Welche Formel passt zur Figur? ($a$, $c$ parallele Seiten, $h$ Höhe, $e$, $f$ Diagonalen) \\ \kreuz{$A = a \cdot h$} \\ \kreuz{$A = \frac{a + c}{2} \cdot h$} \\ \kreuz{$A = \frac{e \cdot f}{2}$}

\trapez[hoehe=h]{5}{3}{2.5}
\teil Trapez, $a = 9$ cm, $c = 5$ cm, $h = 4$ cm – Fläche? A = \leerfeld[cm²]
\teil Trapez, $a = 7{,}5$ cm, $c = 4{,}5$ cm, $h = 3$ cm – Fläche? A = \leerfeld[cm²]
\end{teile}
\end{aufgabe}

\begin{aufgabe}{Trapez – Prüfungsaufgaben}
\begin{teile}
\steil Ein Regalbrett hat die Form eines gleichschenkligen Trapezes mit den parallelen Seiten $90$ cm und $70$ cm; seine Fläche ist $2\,960$ cm². Passen Blumentöpfe mit $35$ cm Durchmesser darauf? Weise es über die Tiefe (Höhe) nach. \hfill (P10 2014 OS)
\steil Eine Terrasse hat die Form eines Trapezes; sie ist $16{,}2$ m² groß, die parallelen Seiten sind $7{,}4$ m und $4{,}6$ m lang. Reicht die Tiefe (Höhe) für eine $3$ m lange Liege? Weise es nach. \hfill (P10 2014 OS)
\end{teile}
\end{aufgabe}

% Anlauf (ohne Prüfkennung)
\begin{aufgabe}{Zusammengesetzte Figuren – Anlauf}
\begin{teile}
\teil In welche Teilflächen zerlegst du die Figur? Benenne sie.

\viereck{(0,0)}{(6,0)}{(6,3)}{(2,3)}
\teil Ein L-förmiges Beet besteht aus zwei Rechtecken: $8$ m × $3$ m und $3$ m × $4$ m. Fläche? A = \leerfeld[m²]
\teil Eine Giebelwand: unten ein Rechteck, $8$ m breit und $3$ m hoch, darauf ein Dreieck mit derselben Grundseite und $2{,}5$ m Höhe. Fläche der Wand? A = \leerfeld[m²]
\end{teile}
\end{aufgabe}

\begin{aufgabe}{Zusammengesetzte Figuren – Prüfungsaufgaben}
\begin{teile}
\teil Eine Eisbahn besteht aus einem Rechteck $24$ m × $14$ m und an beiden kurzen Seiten je einem Halbkreis mit $14$ m Durchmesser. Berechne die Fläche der Eisbahn. Runde auf ganze m². \hfill (P10 2017 OS) \\ A = \leerfeld[m²]
\teil Ein Beet besteht aus einem Rechteck $6{,}5$ m × $3$ m und an beiden kurzen Seiten je einem Halbkreis. Berechne die Fläche des Beets. Runde auf eine Stelle. \hfill (P10 2017 OS) \\ A = \leerfeld[m²]
\end{teile}
\end{aufgabe}
